\par\medskip
\hypertarget{native-076-article058}{}\pdfbookmark[2]{58.}{article-058}
\begin{center}58.\end{center}
Formulae nostrae pro distantiis a planis datis involvunt $v$ et $r$: quoties has quantitates e tempore prius determinare oportet, partem operationum adhuc contrahere, atque sic laborem notabiliter allevare licebit. Derivari enim possunt illae distantiae per formulam persimplicem statim ex anomalia excentrica in ellipsi, vel e quantitate auxiliari $F$ aut $u$ in hyperbola, ita ut computo anomaliae verae radiique vectoris plane non sit opus. Mutatur scilicet expressio $kr\sin(v+K)$

I. \textit{pro ellipsi}, retentis characteribus art. 8, in
\[
ak\cos\varphi\cos K\sin E+ak\sin K(\cos E-e).
\]
Determinando itaque $l$, $L$, $\lambda$ per aequationes
\[
\begin{aligned}
ak\sin K&=l\sin L\\
ak\cos\varphi\cos K&=l\cos L\\
-eak\sin K&=-el\sin L=\lambda,
\end{aligned}
\]
\SourceLeaf{077}expressio nostra transit in $l\sin(E+L)+\lambda$, ubi $l$, $L$, $\lambda$ constantes erunt, quatenus $k$, $K$, $e$ pro constantibus habere licet; sin minus, de illarum mutationibus computandis eadem valebunt, quae in art. praec. monuimus.

Exempli caussa transformationem expressionis pro $x$ in art. 56 inventi apponimus, ubi longitudinem perihelii $=121^0\,17^{\prime}\,34^{\prime\prime}{,}4$, $\varphi=14^0\,13^{\prime}\,31^{\prime\prime}{,}97$, $\log a=0{,}442\,3790$ statuimus. Fit igitur distantia perihelii a nodo ascendente in ecliptica $=308^0\,49^{\prime}\,20^{\prime\prime}{,}7=u-v$; hinc $K=212^0\,36^{\prime}\,56^{\prime\prime}{,}09$. Habemus itaque
\begin{center}
\begin{tabular}{@{}ll@{}}
\makebox[13em][l]{$\log ak$\hfill .\,.\,.\,.\,.\,.\,.} & $0{,}441\,1713$\\[.3em]
\makebox[13em][l]{$\log\sin K$\hfill .\,.\,.\,.\,.} & $9{,}731\,5887_n$\\[.3em]
\makebox[13em][l]{$\log ak\cos\varphi$\hfill .\,.\,.} & $0{,}427\,6456$\\[.3em]
\makebox[13em][l]{$\log\cos K$\hfill .\,.\,.\,.\,.} & $9{,}925\,4698_n$
\end{tabular}
\end{center}
\begin{center}
\begin{tabular}{@{}ll@{}}
\makebox[13em][l]{$\log l\sin L$\hfill .\,.\,.\,.} & $0{,}172\,7600_n$\\[.3em]
\makebox[13em][l]{$\log l\cos L$\hfill .\,.\,.\,.} & $0{,}353\,1154_n$\\[.3em]
\hline
unde $L=$ & $213^0\,25^{\prime}\,51^{\prime\prime}{,}31$\\[.3em]
\makebox[13em][l]{$\log l$\hfill .\,.\,.\,.\,.\,.\,.\,.} & $0{,}431\,6627$\\[.3em]
\makebox[13em][l]{$\log\lambda$\hfill .\,.\,.\,.\,.\,.\,.\,.} & $9{,}563\,2352$\\[.3em]
$\lambda=$ & $+0{,}365\,7929$
\end{tabular}
\end{center}

II. In hyperbola formula $kr\sin(v+K)$ secundum art. 21 transit in
\[
\lambda+\mu\operatorname{tang}F+\nu\operatorname{secans}F,
\]
si statuitur
\[
ebk\sin K=\lambda,\qquad bk\operatorname{tang}\psi\cos K=\mu,\qquad -bk\sin K=\nu;
\]
manifesto eandem expressionem etiam sub formam $\frac{n\sin(F+N)+\nu}{\cos F}$ reducere licet. Si loco ipsius $F$ quantitas auxiliaris $u$ adhibita est, expressio $kr\sin(v+K)$ per art. 21 transibit in $\alpha+\beta u+\frac{\gamma}{u}$, ubi $\alpha$, $\beta$, $\gamma$ determinantur per formulas
\[
\begin{aligned}
\alpha&=\lambda=ebk\sin K\\
\beta&=\tfrac{1}{2}(\nu+\mu)=-\tfrac{1}{2}ebk\sin(K-\psi)\\
\gamma&=\tfrac{1}{2}(\nu-\mu)=-\tfrac{1}{2}ebk\sin(K+\psi).
\end{aligned}
\]

III. In parabola, ubi anomalia vera e tempore immediate derivatur, nihil aliud supererit, nisi ut pro radio vectore valor suus substituatur. Denotando itaque distantiam in perihelio per $q$, expressio $kr\sin(v+K)$ fit $=\frac{qk\sin(v+K)}{\cos\tfrac{1}{2}v^2}\mathord{\cdot}$
